\chapter{Álgebras de conjuntos y medidas}\label{cap:v2-01}
Medir exige decidir primero qué conjuntos se admitirán y después asignarles tamaños compatibles con las operaciones numerables. Una familia de conjuntos no debe confundirse con el espacio que la contiene: usaremos $X$ para el espacio y $\Acal$ para sus conjuntos medibles. Este capítulo desarrolla las lecturas 1 y 2 de Valencia Bucio \cite{valencia2}.

\section{Familias estables bajo operaciones}\label{lectura:v2-01}
\begin{definicion}\index{algebra de conjuntos@álgebra de conjuntos}\index{sigmaálgebra@\(\sigma\)-álgebra}
Un \emph{álgebra} sobre $X$ es una familia $\Acal\subseteq\mathcal P(X)$ que contiene $X$ y es estable bajo complementos relativos a $X$ e intersecciones finitas. Una \emph{$\sigma$-álgebra} es un álgebra estable, además, bajo uniones numerables.
\end{definicion}
Toda álgebra contiene $\varnothing=X\setminus X$. Por De Morgan contiene $A\cup B=(A^c\cap B^c)^c$; por inducción contiene uniones e intersecciones finitas. También contiene $A\setminus B=A\cap B^c$ y $A\mathbin\triangle B=(A\setminus B)\cup(B\setminus A)$. En una $\sigma$-álgebra, De Morgan da igualmente intersecciones numerables. Estas identidades explican cada propiedad: no son nuevas condiciones independientes.

\begin{ejemplo}
Las familias $\{\varnothing,X\}$ y $\mathcal P(X)$ son $\sigma$-álgebras. Los conjuntos finitos o cofinitos forman un álgebra: el complemento intercambia ambas clases; la unión de dos finitos es finita y una unión que contiene un cofinito es cofinita. Si $X$ es infinito, esta familia puede no ser una $\sigma$-álgebra. En $\N$, la unión de los singletons pares es infinita y tiene complemento infinito.

Más generalmente, para un cardinal infinito $\kappa$, los conjuntos de cardinal menor que $\kappa$ o con complemento de cardinal menor que $\kappa$ forman un álgebra. La unión de dos conjuntos de cardinal menor que $\kappa$ sigue teniendo cardinal menor que $\kappa$: el máximo de sus dos cardinales, multiplicado por $2$, permanece menor que $\kappa$. Los conjuntos de cardinal a lo sumo $\kappa$ o con complemento de cardinal a lo sumo $\kappa$ forman una $\sigma$-álgebra, pues una unión numerable de conjuntos de cardinal a lo sumo $\kappa$ tiene cardinal a lo sumo $\kappa\aleph_0=\kappa$. Se usa elección en esta aritmética cardinal.
\end{ejemplo}
\begin{observacion}[Una familia que no es un álgebra]
En la lectura 1, ejemplo 3(iii), se afirma que los intervalos y sus complementos forman un álgebra en $\R$. La afirmación es falsa: $(0,1)\cup(2,3)$ no es un intervalo y su complemento tampoco lo es. Una versión correcta es la familia de \emph{uniones finitas de intervalos}, permitiendo intervalos no acotados, singletons y cualquier combinación de extremos abiertos o cerrados. Una partición finita de la recta por todos los extremos involucrados verifica su estabilidad bajo las operaciones del álgebra.
\end{observacion}

\begin{proposicion}\label{prop:v2-generada}
La intersección de una familia no vacía de álgebras sobre $X$ es un álgebra; lo mismo vale para $\sigma$-álgebras. Por tanto, para cualquier $\mathcal G\subseteq\mathcal P(X)$ existen el álgebra $A(\mathcal G)$ y la $\sigma$-álgebra $\sigma(\mathcal G)$ generadas por $\mathcal G$: son las menores familias de la clase correspondiente que contienen $\mathcal G$.
\end{proposicion}
\begin{proof}
El conjunto $X$ pertenece a cada familia y, por ello, a la intersección. Si $A$ pertenece a todas, también $A^c$ pertenece a todas. Una intersección finita o unión numerable de elementos comunes pertenece a cada familia por su estabilidad. Para construir la familia generada intersectamos todas las familias admisibles que contienen $\mathcal G$. Hay al menos una, $\mathcal P(X)$; la intersección contiene $\mathcal G$ y está contenida en cualquier otra candidata.
\end{proof}
\begin{definicion}\index{Borel, sigmaalgebra de@Borel, sigmaálgebra de}
En un espacio topológico $X$, $\Bcal(X)=\sigma(\tau)$ es la \emph{$\sigma$-álgebra de Borel}. Para espacios medibles $(X,\Acal)$ y $(Y,\Bcal)$ se define
\[\Acal\otimes\Bcal=\sigma\{A\times B:A\in\Acal,\ B\in\Bcal\}.\]
La familia de rectángulos generadores generalmente no es un álgebra.
\end{definicion}

\section{Ejercicios sobre las familias medibles}
\begin{ejercicio}[Lectura 1, ejercicio 1]\label{ej:v2-l01-1}
Sea $\mathcal G$ una familia no vacía estable bajo complementos y uniones binarias. Pruebe que es un álgebra.
\end{ejercicio}

\begin{ejercicio}[Lectura 1, ejercicio 2]\label{ej:v2-l01-2}
Dada una sucesión $(A_n)$ en un álgebra, construya $B_n\subseteq A_n$, disjuntos dos a dos, con $\bigcup_nB_n=\bigcup_nA_n$.
\end{ejercicio}

\begin{ejercicio}[Lectura 1, ejercicio 3]\label{ej:v2-l01-3}
Pruebe que $\Acal|_Y=\{A\cap Y:A\in\Acal\}$ es un álgebra, o una $\sigma$-álgebra, sobre $Y\subseteq X$ si $\Acal$ lo es sobre $X$.
\end{ejercicio}

\begin{ejercicio}[Lectura 1, ejercicio 4]\label{ej:v2-l01-4}
Enliste todas las álgebras sobre un conjunto de cuatro elementos.
\end{ejercicio}

\begin{ejercicio}[Lectura 1, ejercicio 5]\label{ej:v2-l01-5}
Caracterice $A(\mathcal G)$ mediante expresiones finitas de intersecciones y complementos de elementos de una familia no vacía $\mathcal G$.
\end{ejercicio}

\begin{ejercicio}[Lectura 1, ejercicio 6]\label{ej:v2-l01-6}
En un espacio métrico separable, pruebe que las bolas abiertas generan $\Bcal(X)$.
\end{ejercicio}

\begin{ejercicio}[Lectura 1, ejercicio 7]\label{ej:v2-l01-7}
Pruebe que los rayos $(a,\infty)$ generan $\Bcal(\R)$.
\end{ejercicio}

\begin{ejercicio}[Lectura 1, ejercicio 8]\label{ej:v2-l01-8}
Si $X$ tiene al menos dos puntos, pruebe que los rectángulos $A\times B$ no forman un álgebra en $X\times X$. Determine el álgebra que generan.
\end{ejercicio}

\begin{ejercicio}[Lectura 1, ejercicio 9]\label{ej:v2-l01-9}
Si $f:X\to Y$ y $\Bcal$ es una $\sigma$-álgebra sobre $Y$, determine si $\{f^{-1}(B):B\in\Bcal\}$ es una $\sigma$-álgebra.
\end{ejercicio}


\section{Aditividad numerable y continuidad}\label{lectura:v2-02}
En $[0,\infty]$ adoptamos $0\cdot\infty=0$ y $a+\infty=\infty$. Nunca efectuaremos la resta $\infty-\infty$. Una serie de términos no negativos es el supremo de sus sumas parciales, incluso cuando vale infinito.
\begin{definicion}\index{medida}\index{sigmafinitud@\(\sigma\)-finitud}
Una \emph{medida} en $(X,\Acal)$ es una función $\mu:\Acal\to[0,\infty]$ tal que $\mu(\varnothing)=0$ y
\[\mu\left(\bigcup_nA_n\right)=\sum_n\mu(A_n)\quad\text{si los }A_n\text{ son disjuntos dos a dos}.\]
La terna $(X,\Acal,\mu)$ es un espacio de medida. Es finito si $\mu(X)<\infty$; es $\sigma$-finito si $X$ admite una partición numerable en conjuntos medibles de medida finita.
\end{definicion}
\begin{ejemplo}
La medida nula vale cero en todos los conjuntos. La medida de conteo asigna a un conjunto su número de elementos, o infinito si es infinito: en una unión disjunta los elementos se cuentan una sola vez. Para $x_0\in X$ y $\alpha\in[0,\infty]$, la medida puntual $\alpha\delta_{x_0}$ vale $\alpha$ cuando el conjunto contiene $x_0$ y cero en caso contrario. En una familia disjunta, $x_0$ pertenece a lo sumo a un término, lo que prueba la aditividad. Una medida de probabilidad es una medida con $\mu(X)=1$.\index{medida!de conteo}\index{medida!de Dirac}
\end{ejemplo}
\begin{teorema}\label{teo:v2-contmedida}\index{continuidad!de una medida}
Una medida es monótona y numerablemente subaditiva. Si $A_n\uparrow A$, entonces $\mu(A_n)\uparrow\mu(A)$. Si $A\subseteq B$ y $\mu(B)<\infty$, entonces $\mu(B\setminus A)=\mu(B)-\mu(A)$.
\end{teorema}
\begin{proof}
De $B=A\mathbin{\dot\cup}(B\setminus A)$ resulta $\mu(B)=\mu(A)+\mu(B\setminus A)\ge\mu(A)$. Si $\mu(B)$ es finita, ambos sumandos son finitos y se puede restar. Si $A_n\uparrow A$, ponga $D_1=A_1$, $D_n=A_n\setminus A_{n-1}$: son disjuntos y $A_n=\bigcup_{j\le n}D_j$. Por ello $\mu(A)=\sum_j\mu(D_j)=\lim_n\mu(A_n)$, con igualdad también si la serie diverge a infinito. Para una sucesión arbitraria use, en cambio, $D_n=A_n\setminus\bigcup_{j<n}A_j$. Entonces
\[\mu\left(\bigcup_nA_n\right)=\sum_n\mu(D_n)\le\sum_n\mu(A_n).\]
La diferencia con todos los términos anteriores es indispensable: quitar solamente $A_{n-1}$ no produce disjunción en general. Esta es la corrección necesaria en la prueba de subaditividad de la lectura 2.
\end{proof}

\section{Ejercicios sobre medidas y límites de conjuntos}
\begin{ejercicio}[Lectura 2, ejercicio 1]\label{ej:v2-l02-1}
Para $A\in\Acal$ fijo, pruebe que $\lambda(B)=\mu(A\cap B)$ es una medida.
\end{ejercicio}

\begin{ejercicio}[Lectura 2, ejercicio 2]\label{ej:v2-l02-2}
Pruebe que las combinaciones finitas no negativas de medidas son medidas. Pruebe también que $\lambda=\sum_n\alpha_n\mu_n$ es una medida si $\alpha_n\ge0$ y las $\mu_n$ son medidas finitas.
\end{ejercicio}

\begin{ejercicio}[Lectura 2, ejercicio 3]\label{ej:v2-l02-3}
Si $A_n\downarrow A$ y algún $A_N$ tiene medida finita, pruebe que $\mu(A_n)\to\mu(A)$. Muestre la necesidad de la hipótesis de finitud.
\end{ejercicio}

\begin{ejercicio}[Lectura 2, ejercicio 4]\label{ej:v2-l02-4}
Pruebe $\mu(\liminf A_n)\le\liminf\mu(A_n)$, donde $\liminf A_n=\bigcup_m\bigcap_{n\ge m}A_n$.
\end{ejercicio}

\begin{ejercicio}[Lectura 2, ejercicio 5]\label{ej:v2-l02-5}
Si $\mu(\bigcup_nA_n)<\infty$, pruebe $\limsup\mu(A_n)\le\mu(\limsup A_n)$.
\end{ejercicio}

\begin{ejercicio}[Lectura 2, ejercicio 6]\label{ej:v2-l02-6}
Defina $\mu_*(A)=\sup\{\mu(B):B\in\Acal,\ B\subseteq A,\ \mu(B)<\infty\}$. Pruebe que es una medida y que coincide con $\mu$ si esta es $\sigma$-finita.
\end{ejercicio}

\begin{ejercicio}[Lectura 2, ejercicio 7: Borel--Cantelli]\label{ej:v2-l02-7}\index{Borel--Cantelli, lema de}
Si $\sum_n\mu(A_n)<\infty$, pruebe que $\mu(\limsup A_n)=0$.
\end{ejercicio}

\begin{ejercicio}[Lectura 2, ejercicio 8]\label{ej:v2-l02-8}
Pruebe que la $\sigma$-finitud equivale a tener $E_n\uparrow X$ con $\mu(E_n)<\infty$.
\end{ejercicio}
